\documentclass[11pt]{article}
\usepackage{amsmath, amssymb, amsthm}
\usepackage{geometry}
\usepackage{mdframed}
\geometry{margin=1in}

\title{Informal Derivation of Euler's Identity}
\author{David Rosenman}
\date{August 2, 2018}


\begin{document}

\maketitle

\begin{mdframed}
\emph{%
“Who would have thought that $\pi$ which enters as the ratio of
circumference to diameter, $e$, as the natural base for logarithms,
$i$, as the fundamental imaginary unit and 0 and 1 (which we know
all about from infancy) would all be tied together in any way,
not to mention such a simple and compact way? I hope I never
stumble into anything like this formula, for nothing I do after
that in life would have any significance.”
}

\hfill — Ramamurti Shankar
\end{mdframed}

The Nobel Prize-winning physicist Richard Feynman called equation (1), known as Euler's identity, "one of the most remarkable, almost astounding, formulas in all of mathematics."

\begin{equation}
    {e^{i\pi }} + 1 = 0
\end{equation}

Euler's identity can be derived easily from the Taylor series...
\begin{equation}
f\left( x \right) = \sum\limits_{n = 0}^\infty  {\frac{{{f^{\left( n \right)}}\left( 0 \right)}}{{n!}}{x^n}} 
\end{equation}

...as well as defining $e^{ax}$ as
\begin{equation}
    {e^{ax}} = \sum\limits_{n = 0}^\infty  {\frac{{{{\left( {ax} \right)}^n}}}{{n!}}} 
\end{equation}

In equation (2), $f^(n)$ represents the $n$-th derivative of the function $f$, with the zeroth derivative of $f$ being equal to $f$.

The Taylor series of $cos(x)$ can be easily derived from equation (2).

\[
\renewcommand{\arraystretch}{1.25} % try 1.2–1.6 to taste
\begin{array}{l}
{f^0}\left( x \right) = f\left( x \right) = \cos x\\
{f^0}\left( 0 \right) = f\left( 0 \right) = \cos 0 = 1\\
{f^1}\left( x \right) = \frac{{df\left( x \right)}}{{dx}} = \frac{d}{{dx}}\cos x =  - \sin x\\
{f^1}\left( 0 \right) =  - \sin 0 = 0\\
{f^2}\left( x \right) = \frac{{{d^2}f\left( x \right)}}{{dx^2}} =  - \frac{d}{{dx}}\sin x =  - \cos x\\
{f^2}\left( 0 \right) =  - \cos 0 =  - 1\\
{f^3}\left( x \right) =  - \frac{d}{{dx}}\cos x = \sin x\\
{f^3}\left( 0 \right) = \sin 0 = 0\\
{f^4}\left( x \right) = \frac{d}{{dx}}\sin x = \cos x\\
{f^4}\left( 0 \right) = \cos 0 = 1\\
\end{array}
\]

From the equations above, as well at equation (2), it follows that:
\begin{equation}
    \cos x = 1 - \frac{{{x^2}}}{{2!}} + \frac{{{x^4}}}{{4!}} - ... = \sum\limits_{n = 0}^\infty  {\frac{{{{\left( { - 1} \right)}^n}}}{{\left( {2n} \right)!}}}  {x^{2n}}
\end{equation}

The Taylor series of $\sin x$ can also be derived from equation (2).

\[
\renewcommand{\arraystretch}{1.25} % try 1.2–1.6 to taste
\begin{array}{l}
{f^0}\left( x \right) = f\left( x \right) = \sin x\\
{f^0}\left( 0 \right) = \sin 0 = 0\\
{f^1}\left( x \right) = \frac{d}{{dx}}\sin x = \cos x\\
{f^1}\left( 0 \right) = \cos 0 = 1\\
{f^2}\left( x \right) = \frac{d}{{dx}}\cos x =  - \sin x\\
{f^2}\left( 0 \right) =  - \sin 0 = 0\\
{f^3}\left( x \right) =  - \frac{d}{{dx}}\sin x =  - \cos x\\
{f^3}\left( 0 \right) =  - \cos 0 =  - 1\\
{f^4}\left( x \right) =  - \frac{d}{{dx}}\cos x = \sin x\\
{f^4}\left( 0 \right) = \sin 0 = 0\\
{f^5}\left( x \right) = \frac{d}{{dx}}\sin x = \cos x\\
{f^5}\left( 0 \right) = \cos 0 = 1
\end{array}
\]

From the equations above and equation (2), it follows that:

\begin{equation}
    \sin x = x - \frac{{{x^3}}}{{3!}} + \frac{{{x^5}}}{{5!}} - ... = \sum\limits_{n = 0}^\infty  {\frac{{{{\left( { - 1} \right)}^n}}}{{\left( {2n + 1} \right)!}}{x^{2n + 1}}} 
\end{equation}

From equations (3), (4), and (5), if $\theta$ is a real number:

\[
\renewcommand{\arraystretch}{1.25} % try 1.2–1.6 to taste
\begin{array}{l}
e^{i\theta} = \sum\limits_{n=0}^{\infty} \frac{(i\theta)^n}{n!}
= 1 + i\theta + \underbrace{i^2}_{-1}\frac{\theta^2}{2!}
+ \underbrace{i^3}_{-i}\frac{\theta^3}{3!}
+ \underbrace{i^4}_{1}\frac{\theta^4}{4!}
+ \underbrace{i^5}_{i}\frac{\theta^5}{5!} - \ldots\\
e^{i\theta} = 1 + i\theta - \frac{\theta^2}{2!} - \frac{i\theta^3}{3!}
+ \frac{\theta^4}{4!} + \frac{i\theta^5}{5!} - \ldots
\end{array}
\]

\begin{equation}
    e^{i\theta} =
\underbrace{\left(1 - \frac{\theta^2}{2!} + \frac{\theta^4}{4!} - \ldots\right)}_{\cos\theta}
+ i\underbrace{\left(\theta - \frac{\theta^3}{3!} + \frac{\theta^5}{5!} - \ldots\right)}_{\sin\theta}
= \cos\theta + i\sin\theta
\end{equation}

Finally, let $\theta = \pi$. Equation (6) becomes:
\begin{equation*}
    {e^{i\pi }} = \cos \pi  + i\sin \pi  =  - 1
\end{equation*}

And finally:
\begin{equation*}
    \boxed{e^{i\pi} + 1 = 0}
\end{equation*}


\section*{References}
\begin{enumerate}
  \item Ramamurti Shankar, \emph{Principles of Quantum Mechanics}. Springer, 1994.
  \item James Stewart, \emph{Calculus: Early Transcendentals}. Cengage Learning, 2012.
\end{enumerate}


\end{document}
